\chapter{Hardware Microarchitecture}

\section*{Purpose}
This chapter describes the current PL-side organization at a level suitable for integration planning and early verification.

\section{Current Block Structure}
\begin{itemize}[leftmargin=*]
  \item \textbf{AXI-Lite wrapper:} register decode, digest staging, command acceptance, status reporting, and signature readback windows.
  \item \textbf{ML-DSA engine adapter:} stable internal boundary that accepts wrapper-level start and digest inputs and returns busy, done, error, signature length, and signature data.
  \item \textbf{ML-DSA-OSH shim:} project-owned stream adapter that formats the appliance digest input, partitions PoC key material, sequences the imported sign transaction, and reorders signature output for the wrapper-visible buffer.
  \item \textbf{Imported ML-DSA-OSH core path:} vendored third-party sign implementation under `hw/ip/mldsa\_osh/upstream/`.
  \item \textbf{Signature buffer exposure:} wrapper-visible byte window populated from adapter output once an operation completes.
\end{itemize}

\section{Current Adapter Modes}
The adapter supports three modes:
\begin{itemize}[leftmargin=*]
  \item \textbf{STUB mode:} deterministic fake-signature behavior used to validate wrapper sequencing and software integration.
  \item \textbf{CORE\_PLACEHOLDER mode:} deterministic seam-testing mode that preserves the same adapter contract while standing in for downstream core attachment.
  \item \textbf{MLDSA\_OSH mode:} real attachment path that instantiates project-owned shim logic for the imported ML-DSA-OSH sign interface.
\end{itemize}

\section{Clocking and Reset}
The initial implementation assumes a single primary PL clock domain for the wrapper and adapter path. Reset behavior forces the wrapper into an idle, non-busy state with cleared status and an empty signature window. The ML-DSA-OSH shim additionally applies a short internal reset sequence before launching the upstream sign operation.

\section{Current Dataflow}
Software writes the digest into wrapper-visible registers and asserts `start`. The wrapper launches the adapter. In `STUB` and `CORE\_PLACEHOLDER` modes, the adapter returns a deterministic local result. In `MLDSA\_OSH` mode, the shim converts the digest-oriented appliance request into the inspected upstream sign stream, captures the real core output stream, and reorders it into the wrapper-visible signature buffer.

\section{Digest Adaptation}
The appliance boundary remains digest-oriented, while the inspected upstream sign path expects formatted message bytes. The current shim therefore uses a provisional PoC adaptation that emits `0x00 || 0x00 || digest[63:0]` as the formatted message payload and provides `mlen = 64`. This assumption is isolated inside project-owned shim logic and documented explicitly so it can be revisited later without changing the wrapper or software interface.

\section{Signature Buffering}
The deterministic modes still report 128-byte signatures. The real ML-DSA-87 path uses a 4627-byte signature. To keep the wrapper contract stable while supporting the real path, the `SIG\_DATA` window has been compatibly sized for the largest supported result. The `SIG\_LENGTH` register continues to report the actual valid byte count for the selected engine mode.

\section{Adapter Rationale}
The adapter prevents the AXI-Lite wrapper from depending directly on ML-DSA-OSH-specific handshake, key-segment ordering, or output-stream ordering details. This allows the PS-visible contract to remain stable while the downstream engine implementation evolves from deterministic scaffolding to a real hardware core.

\section{Current Verification Constraint}
The imported ML-DSA-OSH sign path mixes Verilog and VHDL sources. The current local portable toolchain is sufficient for wrapper-level regression and for honest testing of the ML-DSA-OSH fallback behavior, but not for full local simulation of the mixed-language imported sign path. That limitation is tooling-related rather than an architectural decision.